% Edited conversion of source pp. 38--39.
\ReportHeading
  {Аналіз п’єзоелектрика з міжфазною тріщиною, заповненою частково електропроникним матеріалом}
  {М.~С. Левченко, В.~В. Лобода}
  {Дніпровський національний університет імені Олеся Гончара}
  {}

Розглянуто задачу плоскої деформації у площині $(x_1,x_3)$ для двох з’єднаних прямокутних областей $x_3>0$ та $x_3<0$, кожна з яких має висоту $l$ і ширину $2l$. На межі поділу $x_3=0$ розташована тріщина $-b<x_1<b$. Обидва матеріали належать до класу симетрії $6mm$, а їхні осі поляризації орієнтовано вздовж $x_3$. На зовнішніх сторонах прямокутників задано нормальні та дотичні напруження й електричне зміщення, незалежні від координати $x_2$.

Тріщина вважається повністю відкритою; її береги вільні від заданих механічних навантажень та електричних зарядів. Заповнювач має діелектричну проникність
\[
  \varepsilon_a=\varepsilon_r\varepsilon_0,
  \qquad
  \varepsilon_0=8.85\times10^{-12}\ \mathrm{C/(V\,m)},
\]
де $\varepsilon_r$ --- відносна діелектрична проникність, а $\varepsilon_0$ --- діелектрична проникність вакууму.

Тріщину моделюють дуже тонким еліпсом, вертикальна вісь якого має порядок $10^{-5}$ від довжини горизонтальної осі. Еліптичну порожнину заповнюють гіпотетичним п’єзоелектричним матеріалом, модулі пружності та п’єзоелектричні сталі якого мають порядок $10^{-6}$ від відповідних характеристик верхнього матеріалу. Його діелектричні проникності, навпаки, задають рівними проникностям реальних матеріалів, які можуть заповнювати тріщину. Така модель з високою точністю відтворює математичний розріз зі скінченною електричною проникністю заповнювача~\cite{piezo:Hao1994}.

Модель досліджено методом скінченних елементів. Для порівняння з аналітичними результатами~\cite{piezo:Govorukha2006} довжину тріщини вважали значно меншою за характерні розміри областей. Розрахункова область складається з трьох ідеально з’єднаних підобластей із різними фізичними характеристиками. Використано восьмивузлові чотирикутні та шестевузлові трикутні скінченні елементи лагранжевого типу зі згущенням сітки біля тріщини, особливо в околі її вершин.

Для різних комбінацій верхнього й нижнього п’єзоелектричних матеріалів та заповнювачів отримано електричний потік через тріщину, стрибки переміщень і електричного потенціалу на її берегах. Значення потоку добре узгоджуються з аналітичною моделлю тріщини зі скінченною електричною проникністю. Якщо діелектрична проникність заповнювача не менша за проникність води, різниця потенціалів на берегах майже зникає для всіх розглянутих комбінацій матеріалів, тобто тріщина стає практично електропроникною.

\textit{Роботу виконано за фінансової підтримки Національного фонду досліджень України, грант №~2025.07/0117.}

\begin{thebibliography}{9}
\bibitem{piezo:Hao1994}
T.~H. Hao and Z.~Y. Shen,
``A new electric boundary condition of electric fracture mechanics and its applications,''
\emph{Engineering Fracture Mechanics}, vol.~47, pp.~793--802, 1994.

\bibitem{piezo:Govorukha2006}
V.~B. Govorukha, V.~V. Loboda, and M.~Kamlah,
``On the influence of the electric permeability on an interface crack in a piezoelectric bimaterial compound,''
\emph{International Journal of Solids and Structures}, vol.~43, pp.~1979--1990, 2006.
\end{thebibliography}
